\documentclass[12pt]{article}
\usepackage{geometry} % Pour passer au format A4
\geometry{hmargin=1cm, vmargin=1cm} % 

% Page et encodage
\usepackage[T1]{fontenc} % Use 8-bit encoding that has 256 glyphs
\usepackage[english,french]{babel} % Français et anglais
\usepackage[utf8]{inputenc} 

\usepackage{lmodern}
\setlength\parindent{0pt}

% Graphiques
\usepackage{graphicx,float,grffile}

% Maths et divers
\usepackage{amsmath,amsfonts,amssymb,amsthm,verbatim}
\usepackage{multicol,enumitem,url,eurosym,gensymb}

% Sections
\usepackage{sectsty} % Allows customizing section commands
\allsectionsfont{\centering \normalfont\scshape}

% Tête et pied de page

\usepackage{fancyhdr} 
\pagestyle{fancyplain} 

\fancyhead{} % No page header
\fancyfoot{}

\renewcommand{\headrulewidth}{0pt} % Remove header underlines
\renewcommand{\footrulewidth}{0pt} % Remove footer underlines

\newcommand{\horrule}[1]{\rule{\linewidth}{#1}} % Create horizontal rule command with 1 argument of height

%----------------------------------------------------------------------------------------
%   Début du document
%----------------------------------------------------------------------------------------

\begin{document}

%----------------------------------------------------------------------------------------
% RE-DEFINITION
%----------------------------------------------------------------------------------------
% MATHS
%-----------

\newtheorem{Definition}{Définition}
\newtheorem{Theorem}{Théorème}
\newtheorem{Proposition}{Propriété}

% MATHS
%-----------
\renewcommand{\labelitemi}{$\bullet$}
\renewcommand{\labelitemii}{$\circ$}
%----------------------------------------------------------------------------------------
%   Titre
%----------------------------------------------------------------------------------------

\setlength{\columnseprule}{1pt}

\horrule{2px}
\section*{Chapitre 3 - Puissances}
\horrule{2px}

\begin{enumerate}
\item[1.] Connaissances
  \begin{itemize}
  \item Opération Puissance positive : $x^n$.
  \item Opération Puissance négative : $x^{-n}$.
  \item Connaître les simplifications avec des puissances.
  \item Connaître l'écriture scientifique.
  \end{itemize}
\item[2.] Compétences
  \begin{itemize}
  \item Calculer avec des puissances.
  \item Comprendre quelques limites de sa calculatrice.
  \end{itemize}
\end{enumerate}

\section*{1 - L'opération puissance}

\begin{multicols}{2}

  \subsection*{Les puissances de 2}

  Les puissances de 2 sont utilisées à la base de l'informatique pour stocker des nombres dans le système de numération binaire.

  \textbf{L'opération puissance est une nouvelle opération}. On multiplie un nombre par lui-même un certain nombre de fois. En ce point, elle ressemble à la multiplication. \textit{On lit deux puissance quatre}. 

  \begin{flalign*}
     2 + 2 + 2 + 2 &= 2 \times 4 \\
    2 \times 2 \times 2 \times 2 &= 2^4
  \end{flalign*}

  \underline{Liste des première puissances de 2 :} \\

  \begin{itemize}
  \item $2^0 = 1 $ // Par définition
  \item $2^1 = 2$
  \item $2^2 = 2 \times 2 = 4$
  \item $2^3 = 2 \times 2 \times 2 = 8$
  \item $2^4 = 2 \times 2 \times 2 \times 2 = 16 $
  \item $2^5 = 2 \times 2 \times 2 \times 2 \times 2 = 32$
  \item $2^6 = 2 \times ... \times 2 = 64$
  \item $2^7 = 2 \times ... \times 2 = 128$
  \item $2^8 = 2 \times ... \times 2 = 256$
  \item $2^9 = 2 \times ... \times 2 = 512$
  \item $2^{10} = 2 \times ... \times 2 = 1024$
  \end{itemize}

  \subsection*{Les puissances de 2 négatives}

  L'opération puissance fonctionne aussi avec des \textbf{exposants négatifs}. On divise alors un certain nombre de fois par le nombre de départ. Elle fait appel aux fractions.

  \begin{flalign*}
        -2 - 2 - 2 - 2 &= 2 \times -4 \\
    1 \div 2 \div 2 \div 2 \div 2 &= 2^{-4}
  \end{flalign*}

  \underline{Liste des premières puissances de 2 négatives :} \\

  \begin{itemize}
  \item $2^0 = 1 $ 
  \item $2^{-1} = \dfrac{1}{2^{1}} = \dfrac{1}{2} $
  \item $2^{-2} = \dfrac{1}{2^{2}} = \dfrac{1}{4} $
  \item $2^{-3} = \dfrac{1}{2^{3}} = \dfrac{1}{8} $
  \item $2^{-4} = \dfrac{1}{2^{4}} = \dfrac{1}{16} $
  \item $2^{-5} = \dfrac{1}{2^{5}} = \dfrac{1}{32} $
  \item $2^{-6} = \dfrac{1}{2^{6}} = \dfrac{1}{64} $
  \end{itemize}

\end{multicols}

\newpage

\section*{2 - La notation scientifique}

\subsection*{Les grands nombres}

\begin{multicols}{2}

  L'opération puissance permet d'obtenir rapidement de très grands nombres.

  $$2^{64} = 18 \, 446 \, 744 \, 073 \, 709 \, 551 \, 616$$

  \paragraph{Problème :}

  \begin{itemize}
  \item Ce nombre est difficile à lire.
  \item Ce nombre est long à représenter.
  \item Ce nombre est source d'erreur à écrit.
  \end{itemize}

  Les calculatrices présente les résultats autrement : \textbf{la forme scientifique.}

  $$2^{64} = 1,844 \, 674 \, 407 \times 10^{19}$$

  \underline{Pour écrire un nombre sous forme scientifique :} \\
  \begin{itemize}
  \item on place la virgule après le premier chiffre puis on multiplie par une puissance de 10. \\
  \end{itemize}

  La puissance de 10 s'appelle \textbf{l'ordre de grandeur.}\\
  \textit{On le lit : un virgule quatre-vingt-quatre fois dix puissances dix-neuf.}

  \paragraph{Conclusion :}

  \begin{itemize}
  \item La représentation de ce nombre tient à l'écran.
  \item Ce nombre est facile à lire et à écrire.
  \end{itemize}

  \subsubsection*{Les puissances de 10}

  \textbf{La notation scientifique} utilise les puissances de dix. Pour connaître la puissance de 10 correspondante, on compte le nombre de chiffre entre le premier chiffre et le chiffre des unités. \\ 

  \underline{Liste des premières puissances de 10 :} \\

  \begin{itemize}
  \item $10^0 = 1 $ // Par définition 
  \item $10^1 = 10$ // On multiplie par dix.
  \item $10^2 = 100$ // On multiplie par cent.
  \item $10^3 = 1 \, 000$ // On multiplie par mille.
  \item $10^4 = 10 \, 000$ // On multiplie par dix mille.
  \item $10^5 = 100 \, 000$ // On multiplie par cent mille.
  \item $10^6 = 1 \, 000\, 000$ // On multiplie par un million.
  \item $10^9 = 1 \, 000\, 000 \, 000$ 
  \item $10^{12} =1 \, 000 \, 000 \, 000 \, 000$
  \end{itemize}

\end{multicols}

\subsection*{Les petits nombres}

\begin{multicols}{2}

  On rencontre les mêmes problèmes avec les petits nombres.

  $$2^{-10} = 0, 000 \, 976\, 562\, 5$$

  Pour les mêmes raisons, les calculatrices vont aussi les représenter sous \textbf{forme scientifique}.

  $$2^{-10} = 9, 7654 \, 625 \times 10^{-4}$$

  \textit{On le lit : neuf virgule soixante seize fois dix puissances moins quatre.}

  \subsubsection*{Les puissances de 10 négatives}

  \textbf{La notation scientifique} utilise les puissances de dix négatives pour représenter les petits nombres. Afin de connaître la puissance de 10 négative, on compte le nombre de 0 en partant du premier. \\

  \underline{Liste des premières puissances de 10 négatives:} \\

  \begin{itemize}
  \item $10^0 = 1 $ // Par définition 
  \item $10^{-1} = 0,1$ // On multiplie par un dixième.
  \item $10^{-2} = 0,01$ // On multiplie par centième.
  \item $10^{-3} = 0,001$ // On multiplie par un millième.
  \item $10^{-4} = 0,000 \, 1$ 
  \item $10^{-5} = 0,000 \, 01$ 
  \item $10^{-6} = 0,000 \, 001$ .
  \item $10^{-9} = 0,000 \, 000 \, 000 \, 1$ 
  \end{itemize}
\end{multicols}

\newpage

\section*{3 - Trucs et astuces}

Il existe \textbf{quatre règles} pour simplifier \textit{à la main} les puissances dans les calculs. Il faudra être prudent avec les puissances négatives et faire attention aux signes. Le but n'est pas de calculer mais de proposer une écriture plus simple.

\setlength{\columnseprule}{0pt}

\begin{multicols}{2}

  \paragraph{Règle 1 : } $ 2^{10} \times 2^{12} = 2^{10 + 12} = 2^{22} $

  \paragraph{Règle 2 : } $ \dfrac{14^{25}}{14^{10}} = 14^{25 - 10} = 14^{15} $

  \paragraph{Règle 3 : } $ 5^{19} \times 7^{19} = (5 \times 7)^{19} = 35^{19} $

  \paragraph{Règle 4 : } $ (10^4)^9 = 10^{4 \times 9} = 10^36 $


\end{multicols}

\section*{4 - La science et le numérique}

Très rapidement, ce chapitre permet d'atteindre les limites de notre calculatrice. \\

\textbf{La limite de puissance} de la calculatrice est atteinte avec un nombre assez grand. Il est possible de faire des calculs avec $10^{99}$. $10^{100}$ donne une erreur. (\textit{On parle d'\textsc{Overflow}.}) \\

$\dfrac{10^{100}}{10^{90}}$ renvoie une erreur à la place de $10^{10}$.\\ 

\textbf{La limite de précision} de la calculatrice est plus rapidement atteinte. La calculatrice fait alors des erreurs de calcul en mathématiques... Pourquoi n'est-ce pas grave ?

\subsubsection*{Très grand + Normal}

\begin{multicols}{2}

  \begin{figure}[H]
        \centering
        \includegraphics[width=0.8\linewidth]{4x3-puissances/sources/terre.png}
  \end{figure}

  \begin{flalign*}
        10^{16} - 10^{16} + 1.6 &= 1.6 \\
        10^{16} + 1,2 - 10^{16} &= 0 \text{\textbf{ !!!}}
  \end{flalign*}

  Pour nos calculatrices : $10^{16} + 1,6 = 10^{16}$. \\
  Scientifiquement, l'\textbf{approximation} est acceptable. \\
  La distance Terre-Soleil est la même que l'on mesure de notre tête ou de nos pieds.

\end{multicols}

\subsubsection*{Normal + Très petit}

\begin{multicols}{2}

  \begin{figure}[H]
        \centering
        \includegraphics[width=0.8\linewidth]{4x3-puissances/sources/papier.png}
  \end{figure}

  \begin{flalign*}
        12,5 - 12,5 + 10^{-16} &= 10^{-16} \\
        12,5 + 10^{-16} - 12,5 &= 0 \text{\textbf{ !!!}}
  \end{flalign*}

  Pour nos calculatrices : $10^{-16} + 1.6 = 1,6$.\\
  Scientifiquement, l'\textbf{approximation} est acceptable.\\
  Je mesure la même taille si je me mesure du sol ou si je marche sur une feuille de papier.
\end{multicols}

\textbf{Il convient de réfléchir au sens dès qu'on additionne des nombres avec des ordres de grandeurs très différents.}

\end{document}
